\subsection{Universality Criterions}

Given the difficulty of the proof of Theorem \ref{thm:grid-universal}, we now turn our attention to the development of other techniques to determine the universality of a family of resource states. We start by recalling some useful facts regarding regular quantum circuit universality.

\begin{fact}
\label{fact:1-qubit-and-entangling-universal}
Any 1-qubit universal gate set and a 2 qubit entangling gate is universal \cite{1qubit-entangling-universal}.
\end{fact}

\begin{proof}
The proof in \cite{1qubit-entangling-universal} turns out to be quite technical, involving many tricks from Lie Algebra. It surmounts to proving the following restatement of the theorem: Let $V$ be a 2-qubit gate. The following are equivalent:
\begin{enumerate}
    \item The collection of all 1-qubit gates $\cA$ together with $V$ is QQ-universal.
    \item The collection of all 1-qubit gates $\cA$ together with $V$ is exactly QQ-universal.
    \item $V$ is not primitive, or \term{imprimitive}.
\end{enumerate}
Here, say $V$ is \term{primitive} iff $V = S \otimes T$ or $V = (S \otimes T)P$ for some $1$-qubit gates $S$ and $T$, with $P$ being the swap gate.
\end{proof}

\begin{theorem}
\label{thm:any-universal-need-entanglement}
Any universal gate set must include at least one multi-qubit entangling gate, as it is impossible to entangle using only single-qubit gates.
\end{theorem}

\begin{proof}
This follows from Fact \ref{fact:1-qubit-and-entangling-universal}.
\end{proof}

\begin{theorem}
\label{thm:line-not-universal}
Let \term{$\mathsf{Line}$} denote the family of graphs made out of line graphs. Then $\mathsf{Line}$ is not a universal resource.
\end{theorem}

\begin{proof}
This follows from Theorem \ref{thm:any-universal-need-entanglement}.
\end{proof}

We now turn to universality criteria specifically for MBQC. We start by addressing why we only consider CQ-universality rather than QQ-universality.

\begin{theorem}
\label{thm:cq-implies-qq}
This is Observation 2 of \cite{MBQC-universal}. Every CQ-universal resource can be made QQ-universal by allowing quantum inputs to be coupled into the resource state using the same coupling method, i.e. CZ gates.
\end{theorem}

\begin{definition}
Let \term{$\LO$} denote local operations, i.e. any single qubit unitary and measurement. Let \term{$\LC$} denote local Clifford operations, i.e. any ZX diagram with spiders of arity at most $2$ and phases multiples of $\tfrac{\pi}{2}$. Let \term{$\CC$} denote \term{classical communication}, which represents the measurement correction and Pauli error tracking done on-the-fly in MBQC.

States $\ket\psi$ and $\ket\phi$ both on $n$ qubits are called \term{$\LU$-equivalent} if there is a unitary $U = \otimes_{i=1}^n U_i$, where every $U_i$ is a local unitary, such that $\ket\psi = U \ket\phi$. In this case, write $\ket\psi =_\LU \ket\phi$. Now, if every $U_i$ is made up of only local Clifford gates, say they are \term{$\LC$-equivalent} and write $\ket\psi =_\LC \ket\phi$.
\end{definition}

\begin{theorem}
This is the main result of \cite{stab-lc-equiv-graph}. Graph states $\ket{G_1}$ and $\ket{G_2}$ are $\LC$-equivalent iff there is a series of local complementations that transform one graph to another. 
\end{theorem}

\begin{theorem}
This is the main result of \cite{lu-lc-false}. There are graphs that are $\LU$-equivalent but not $\LC$-equivalent. On the other hand, there are some types of graphs which are $\LU$-equivalent iff they are $\LC$-equivalent \cite{lu-equal-lc-example-1,lu-equal-lc-example-2}.
\end{theorem}

\begin{theorem}
A family of states $\Phi$ is a universal resource for MBQC if for all $n$-qubit states $\ket\psi$, $\exists \ket\phi \in \Phi$ with $m \geq n$ qubits, such that the transformation $\ket\phi \to \ket\psi \ket{+^{m-n}}$ is possible deterministically with $\LO$ and $CC$. We denote this by $\ket\phi \geq_{\LOCC} \ket\psi$. 
\end{theorem}

\begin{proof}
This follows almost immediately from the definition. Suppose $\Phi$ is a universal resource. Then the above is certainly true, since a measurement pattern is made up of only $\LO$ and $\CC$. Now suppose the above holds. It is clear that the resource state is universal. Thus the above is a sufficient and necessary condition.
\end{proof}

\begin{theorem}
\label{thm:measure-for-universal}
[From \cite{MBQC-measures}] Let $\Psi$ be a family of resource states and $\ket\phi, \ket{\phi'}$ arbitrary states with the same number of qubits. Let $E: \ket\phi \to \R$ be a function that satisfies the following properties:
\begin{enumerate}
    \item $E(\ket\phi) \geq E(\ket{\phi'})$ whenever $\ket\phi \geq_\LOCC \ket{\phi'}$;
    \item $\sup_{\ket\phi} E(\ket\phi) > \sup_{\ket\psi \in \Psi} E(\ket\psi)$. 
\end{enumerate}
Then $\Psi$ cannot be a universal resource.
\end{theorem}

\begin{proof}
The first condition states that $\LOCC$ cannot increase this measure, while the second says that the supremum value of $E$, which can be $\infty$, is not reached among all $\ket\psi \in \Psi$. If both were true, then clearly universality is not possible.
\end{proof}

The above theorem is very general and consequently very powerful, as we will soon see. However, note that satisfying the converse of the second condition for some measure $E$ is not sufficient to guarantee universality. For example, the family of \emph{toric code states} have unbounded entanglement width but are not universal \cite{toric-not-universal}.

\begin{theorem}
\label{thm:star-not-universal}
Consider the family of all star graphs $\Star$. A \term{star graph} on $n$ vertices is a tree where one vertex has degree $n-1$ and all other vertices have degree $1$. Then $\Star$ is not a universal resource.
\end{theorem}

We choose to prove a more powerful result. 

\begin{lemma}
Every tree has rank width at most 1.
\end{lemma}

\begin{proof}
Let $G=(V,E)$ be a tree. We will construct a rank decomposition where every cut has cut-rank at most $1$. Root the tree $G$ at some vertex $r$. For each vertex $v \neq r$, let $G_v$ denote the subtree consisting of $v$ and all of its descendants.

Let $A_v := V(G_v)$. Since $G$ is a tree, there is exactly one edge between $A_v$ and $V \setminus A_v$, namely the edge from $v$ to its parent Therefore, the cut matrix $M_{A_v}$ has exactly one nonzero entry. Hence
\[
    \rho_G(A_v)
    =
    \operatorname{rank}_{\mathbb F_2} M_{A_v}
    \leq 1.
\]
Now choose a binary decomposition tree $T$ that recursively separates rooted subtrees of $G$. Every cut appearing in this decomposition is of the form $A_v \mid V \setminus A_v$ or a union of such subtrees attached to a common vertex. In either case, because $G$ is a tree, the edges crossing the cut are incident to a single vertex on one side of the cut. Thus the corresponding bipartite adjacency matrix has all nonzero entries contained in a single row or a single column, so its rank over $\mathbb F_2$ is at most $1$. Hence this rank decomposition has width at most $1$. Therefore $\operatorname{rwd}(G) \leq 1$.
\end{proof}

\begin{definition}
Let $\Phi$ and $\Psi$ be families of resource states. We say that $\Phi$ and
$\Psi$ are \term{$\LU$-equivalent}, and write $\Phi =_\LU \Psi$, if the following holds:
\begin{enumerate}
    \item For every $\ket{\phi}\in\Phi$, there exists $\ket{\psi}\in\Psi$ such
    that $\ket{\phi}=_{\LU}\ket{\psi}$, and
    \item for every $\ket{\psi}\in\Psi$, there exists $\ket{\phi}\in\Phi$ such
    that $\ket{\psi}=_{\LU}\ket{\phi}$.
\end{enumerate}
Define \term{$\LC$-equivalent} correspondingly.
\end{definition}

\begin{lemma}
\label{lem:lu-equivalent-families-universality}
Let $\Phi$ and $\Psi$ be $\LU$-equivalent families of resource states. Then $\Phi$ is a universal resource for MBQC if and only if $\Psi$ is a universal resource for MBQC.
\end{lemma}

\begin{proof}
Suppose $\Phi$ is universal, and let $\ket{\rho}$ be any $n$-qubit target state.
By universality, there exists $\ket{\phi}\in\Phi$ on $m\ge n$ qubits such that
$
    \ket{\phi}
    \geq_{\LOCC}
    \ket{\rho}
$
Since $\Phi =_{\LU} \Psi$, there exists $\ket{\psi}\in\Psi$ such that
$\ket{\psi} =_{\LU}\ket{\phi}$. Local unitaries are reversible $\LOCC$
operations, so
$
    \ket{\psi}\geq_{\LOCC}\ket{\phi}.
$
Composing $\LOCC$ transformations gives
$
    \ket{\psi}
    \geq_{\LOCC}
    \ket{\rho}
$
Thus $\Psi$ is universal. The reverse implication is identical.
\end{proof}